\section{The gamma distribution}

To define the gamma distribution, first we define the
$\textbf{Euler gamma function}$ (also called the
$\textbf{generalized factorial function}$), an
important special function of mathematical analysis,
denoted $\Gamma(t)$, by

\[
\Gamma(t) = \int_0^\infty e^{-x} x^{t-1}\,dt \qquad (t>0).
\]

\noindent $\textbf{Properties of the gamma function:}$

\begin{enumerate}[label=\arabic*.]
\item $\Gamma(n)=(n-1)!$ for integer $n\ge1$.

\item $\Gamma(t+1) = t\, \Gamma(t)$ for all $t>0$.

\item $\Gamma(1/2)=\sqrt{\pi}$.
\end{enumerate}

\noindent Next, we say that $X$ has the gamma distribution with
parameters $\alpha, \lambda>0$, and denote
$X\sim \operatorname{Gamma}(\lambda)$, if $X$ has
density function

\[
f_X(x) = \frac{\lambda^\alpha}{\Gamma(\alpha)} e^{-\lambda x} x^{\alpha-1} \qquad (x>0).
\]

\noindent $\textbf{Properties of the gamma distribution:}$

\begin{enumerate}[label=\arabic*.]
\item $\operatorname{Exp}(\lambda) \stackrel{d}{=}
\operatorname{Gamma}(1,\lambda)$.

\item The parameter $\lambda$ has the role of a scale
parameter in the same sense as for the exponential
distribution: for $c>0$ we have
$c \operatorname{Gamma}(\alpha,\lambda)
\stackrel{d}{=} \operatorname{Gamma}(\alpha,
\lambda/c)$.

\item $\operatorname{Gamma}(\alpha,\lambda) \boxplus
\operatorname{Gamma}(\beta,\lambda) =
\operatorname{Gamma}(\alpha+\beta,\lambda)$.

\item $\operatorname{Gamma}(\alpha,\lambda) =
\mathop{\boxplus}\limits_{k=1}^n
\operatorname{Exp}(\lambda)$.

\item If $X\sim \operatorname{Gamma}(\alpha,\lambda)$ then
$\mathbf{E} X = \frac{\alpha}{\lambda}$,
$\mathbf{V}(X) = \frac{\alpha}{\lambda^2}$.
\end{enumerate}
